% SRC: DESIGN.md (stages 1-9); ADR-003; DDD-001; ADR-002
\section{Object pipeline: design and method}

It consumes a Vitrine run read-only and emits \texttt{objects/objects.json}
(schema \texttt{vitrine/object/1}) plus per-object GLBs and a composed scene. Nine stages
run under one evidence-vs-proposal discipline.

\textbf{Identify} runs GroundingDINO text prompts into SAM3.1 promptable segmentation per
frame, degrading to a fixed concept list rather than blocking. \textbf{Rectify} enforces an
explicit image-domain contract: the sidecar persists its own undistorted images and adjusted
intrinsics (because UoS never persists undistorted images), replicating \texttt{undistort.py}
maths under golden tests that require pixel-identical rectification for the same camera. Every
artefact records which pixel domain (distorted or rectified) it lives in.

\textbf{Carve} is the stage where correctness matters most. A 2D mask votes for 3D Gaussians only
through per-pixel front-depth-band gating (foreground / background / \emph{abstain}), with
min-view thresholds, 3D connected-component cleanup and robust p5--p95 extents. The naive
all-along-ray voting of the earlier DreamLab \texttt{mask\_projector} was rejected in adversarial
review because it labels occluded Gaussians behind the object; carve is golden-tested against the
gallery ground-truth object PLYs. \textbf{Seed} scores crops and views by silhouette area,
sharpness, centrality and edge clearance: Lane~A takes the single best matted crop, Lane~B the
top 8--16 diverse views. Splat renders are never used as conditioning.

\textbf{Generate} runs two lanes under one harness (\S\ref{sec:licence}). \textbf{Select}
renders every candidate from all recovered real cameras and scores silhouette IoU, masked
perceptual similarity, depth consistency on observed regions and mesh sanity, reporting observed
and withheld surfaces separately. \textbf{Place} is an owner hard requirement: it initialises
translation and scale from the carve centroid and extents, then refines full SE(3)+scale by
silhouette and photometric alignment against all real masks and frames. It is symmetry-aware and
fits scene-up (EXIF gravity, support-plane fit, or the LiDAR GLB cross-check on
nested-cinema-04) rather than assuming $+Y$. The composed scene rendered from real cameras is the
acceptance proof, and the manifest records per-axis observability, residuals and confidence.

\textbf{Finish} is the second hard requirement: UV unwrap (xatlas/Blender), then projective bake
of real rectified frames onto observed UV islands with per-texel best-view selection by angle,
sharpness and occlusion, and seam blending. Generative fill touches only unseen texels, recorded
in a per-texel provenance mask, followed by albedo delighting. Real frames constrain; generated
pixels propose, always flagged. \textbf{Manifest} writes one lineage record per object: sha256 of
every asset \emph{and} every checkpoint (repo, commit, hash, licence class), stage scores, coverage
estimate, orientation observability and VRAM peaks. No stage may silently substitute generated data
where evidence exists; a missing stage output degrades the record's status field, never fakes it.

\begin{figure}[h]
\centering
% SRC: seed crops runs/nc04-object-pass1/seed/ + Blender renders of runs/nc04-generate-pass1/*/ (scripts/render_turntable.py)
\includegraphics[width=0.86\columnwidth]{renders/assets_pairs.png}
\caption{From one photograph to an asset: the evidence-scored seed crop (left) and the Lane-A
TRELLIS.2 reconstruction rendered from a fresh view (right), for radio, table and speaker. The
generated geometry carries the crop's colour and structure; unobserved faces are inferred.}
\label{fig:assetpairs}
\end{figure}
